\documentclass{article}
\usepackage{graphicx} % Required for inserting images

\title{Freelancing}
\author{Riccardo Parrino}
\date{May 2025}

\begin{document}

\maketitle

\section{Comunicazione}

\subsection{Comunicazione interna}
Rivolta ai dipendenti e ai collaboratori. Ha come obiettivi aumentare la motivazione, la condivisione della cultura aziendale, il coordinamento interno e la gestione del cambiamento. Viene fatto attraverso newsletter, intranet, riunioni, formazione ed eventi interni.

\subsection{Comunicazione esterna}
Rivolta al pubblico esterno all'azienda. Ha come obiettivi promuovere l'immagine aziendale, costruire relazioni, migliorare la reputazione.

Si articola in comunicazione istituzionale, comunicazione commerciale/di marketing, comunicazinoe finanziaria e comunicazione di crisi.

La comunicazione istituzionale cura l'indentit\'a e l'immagine dell'azienda,  ha come target istituzione, investitori ed opinione pubblica. Usa strumenti com e le relazioni pubbliche, gli eventi, i comunicati stampa ed i rapporti di sostenibilit\'a.

Per comunicazione istituzionale, pi\'u precisamente, si intende la comunicazione realizzata in moda organizzato da un'istituzione o dai suoi rappresentanti e diretta alle persone e ai gruppi dell'ambiente sociale in cui svolge la sua attivit\'a. Ha come obiettivo stabilire relazioni di qualit\'a tra l'isituzione e il pubblico con cui si relaziona, per conseguire notoriet\'a sociale e immagine pubblica adeguate ai fini e alle attivit\'a dell'istituzione stessa (Comunicazione pubblica). Nell'ambito pubblicitario, il termine indica la comunicazione volta a promuovere l'immagine di un organismo, commerciale e non, inteso nel suo insieme (corporate image).

La comunicazione commerciale/di marketing mira a promuovere i prodotti o servizi ha come target i clienti attuali e potenziali e usa strumenti come pubblicit\'a , promozioni, direct marketing, digital marketing e social media.

La comunicazione finanziaria \'e rivolta agli investitori, analisti e azionisti, ha contenuti come i risultati economici, bilanci e operazioni finanziarie. 

Infine la comunicazione di crisi mira a comunicaer le situazioni critiche che possono danneggiare l'immagine o la reputazione dell'impresa ed ha come obiettivi la trasparenza, il contenimento dei danni ed il recupero della fiducia.

\subsection{Marketing}
Una visione di insieme del marketing:

\begin{itemize}
    \item Markeing planning and strategy, Business mission and strategy, marketing audit
    \item Marketing objective, strategy, actions
    \item Marketing Evaluation of performance, business performance, metrics
    \item The scope of marketing
\end{itemize}

Understanding Customer Behaviour:
\begin{itemize}
    \item The dimensions of customer behaviour: Who buys, how they buy, what are the choice criteria
    \item  Influences on consumer behaviour, on organizational buying behaviour
\end{itemize}

Marketing Research and Customer Insights
\begin{itemize}
    \item Market research
    \item The role of customer insight
    \item Marketing information systems, market intelligence
\end{itemize}

Market Segmentation, Targeting and Positioning
\begin{itemize}
    \item Segmenting consumer markets, criteria
    \item Segmenting organizational markets, criteria
    \item Target marketing, Positioning, Repositioning
\end{itemize}

Creating Customer Value
\begin{itemize}
    \item Value through Products and Brands
    \item Value through Services, Relationships and Experiences
    \item Value through Pricing
\end{itemize}

Delivering and Managing Customer Value
\begin{itemize}
    \item Delivering Customer Value
    \item Mass Communications Techniques
    \item Direct communications techniques
    \item Digital Marketing
\end{itemize}

\subsection{The Nature of Marketing}
Marketing can be defined as "the achievement of corporate goals through meeting and exceeding customer needs better than the competition". Three conditions must be met before the marketing concept can be applied. First, company activities should be focused on providing customer satisfaction rather than, for example, simply producing products. 

Second,, the achievement of customer satisfaction relies on integrated effort. The responsibility for the implementation of the concept lies not just within the marketing department but should run right through production, finance, research and development, engineering and other departments. Marketing is the responsibility of everyone in the organization.

In summary, companies can be viewed as being either inward looking or outward looking. IN the former, the focus is on making things or providing services  but with significant attention being paid to the efficiency with which internal operatinos and processes are conducted. Companies that build strategy from the outside start with the customer and work backwards from an understanding of what customers tryly value. The difference in emphasis is subtle but very important. By mantainng an outside-in focus, companies can undersand what customers value and how to consistently innovate new sources of value and how to consistently innovate new sources of value that keep bringing them back. Doing so efficiently, ensures that value is created and delivered at a profit to the company - the ultimate goal of marketing. 

Customer value = perceived benefits - perceived sacrifice

Perceived benefits can be derived from the product, the associated service and the image og the company. Conveying benefits is a critical marketing task and is central to positioning and branding. 

Perceived sacrifice is the total cost associated with buying the product. This consists not just of monetary costs, but also the time and energy involved in the purchase. For example, with hotels, good location can reduce the time and energy required to find a suitable place to stay, but marketers need to be aware of another critical sacrifice: the potential psychological cost of not making the right decision. Therefore, hotels like the Marriot or restaurants like McDonald's aim for consistency so that customers can be confident of what they will receive when they visit these service providers. 

A further key to marketing success is to ensure tha the value offered ecveeds that of competitors. Once a product has been purchased, customer satisfaciont depends on its perceived per

Companies need to avoid of setting too high expectations since this can lead to dissatisfaction if performance falls short of expectations. 

Kano model helps to separate characteristics that cause dissatisfaction, satisfaction and delight. 

The three elements detected by the Kano model are "must be", "more is better" and "delighters".

First of all the "must be" component: Those caracteristics recognized as "must bes" are expected and thus taken for granted. For example, commuters expect planes or trains to depart on time and for schedules to be maintained. Lack of these characteristics causes annoyance but their presence only brings dissatisfaction up to a neutral level. "More is better" characteristics can take satisfaction past neutral and into the positive satisfaction range. For example, no response to a telephone call can cause dissatisfaction, but a fast response may cause positive satisfaction. The usability of search results is an example of more is better and has become a key defferentiating factor in the search engine industry, which has allowed Google to becaome the dominant player. "Delighters" are the unexpected characteristics that surprise the customer. Their absence does not cause dissatisfaction, but their presence delights the customer. For exmaple, tourists wha have found that a holiday destination has exceeded their expectations throuhg the quality of customer service that they have received will often be delighted and are likely to recommend the destination to friends and colleagues. 

\subsubsection*{Four forms of customer value}
Modern organizations offer an innnumerable variety of products and services, but four core forms of customer value have been identified as follows: 

\begin{itemize}
    \item Price value
    \item Performance value
    \item Emotional value
    \item Relational value
\end{itemize}

Price value

Performance value

Emotional value: one of the big challenges facing the modern firm is to find effective ways to differentiate products based on performance elements. This helps to explain why some consumers will pay huge premiums for luxury brands (Chanel, Hermes) and why other will queue for hours to be the first among their peers to own certain products (iPad, Kate Moss clothing)

Relation value concers with the quality of service received by the customer. This presents a particular opportunity in the case of servicec businesses such as a restaurant meal or business taxation services which are not easy to evaluate in advnce of purchase.

The challenge for organizations then is to try to become a value leader on one of these four dimesions (Ryanair for price value leader in aviation, Louis Vuitton emotional value leader in luxury fashion goods). This is beacouse they have a clearly defined customer value proposition or unique selling point (USP) which is a reason why customers return to them again and again. 

\subsection{Customer lifecycle management}
\begin{itemize}
    \item Awareness
    \item Consideration
    \item Purchase
    \item Retention
    \item Advocacy
\end{itemize}

\subsection{Cosa fanno le agenzie di marketing e comunicazione}
Le agenzie di marketing aiutano aziende, professionisti o enti a promuovere prodotti, servizi o brand, attraverso strategie e strumenti mirati. Offrono una gamma di servizi che vanno dall'analisi del mercato alla comunicazione pubblicitaria, dal digital marketing  al branding. 

In sintesi, studiano il mercato e i clienti per capire come e dove comunicare, creano strategie per far crescere la visibilit\'a, le vendite e la reputazione, gestiscono campagne pubblicitarie online ed offline, supportano la trasformazione digitale delle aziende. 

I principali servizi offerti dalle agenzie di marketing riguardano la strategia di marketing, il digital marketing, content marketing, branding e identit\'a visiva, grafica e creativit\'a, gestione siti web e e-commerce, pubblicit\'a e campgne media, analisi e reportistica.

\begin{itemize}
    \item Strategia di marketing: analisi SWOT, Studio del target e del posizionamento, piano marketing (tradizionale o digitale)
\end{itemize}

Per piano marketing si intende un documento strategico che definisce obiettivi, azioni, strumenti, tempi e risorse necessarie per promuovere un prodotto, un servizio o un brand, con lo scopo di raggiungere specifici risultati di emrcato (come aumentare le vendite, acquisire clienti, migliorare la notoriet\'a  del marchio).

Un marketing plan solitamente prevede un:

\begin{itemize}
    \item Analisi di contensto: analisi SWOT, del mercato e della concorrenza, analisi del target (chi sono i clienti), posizionamento attuale del brand/prodotto
    \item Marketing Objectives
    \item Stategie di marketing: scelte strategiche su come raggiungere gli obiettivi
    \item Le 4P (Marketing Mix): Prodotto (caratteristiche, packaging, assortimento), Prezzo (politica di pricing, sconti, promozioni), Punto vendita(canali di vendita, e-commerce, retail), promozione (pubblicit\'a, social media, PR, eventi, ecc.)
    \item Piano operativo (azione concreta): Calendario delle attivit\'a (cosa fare, quando, come) , Assegnazione delle responsabilit\'a, Budget e risorse
    \item Monitoraggio e controllo: indicatori chiave di performance (KPI), Analisi dei risultati, Eventuali correzioni di rotta
\end{itemize}

Un esempio molto semplice di marketing plan:
\begin{itemize}
    \item Obiettivo: aumentare le vendite del 30\% nei weekend
    \item Strategia: promuovere i prodotti su Instagram + sconti per chi ordina online
    \item Azioni: 3 post a settimana, stories quotidiane, promozione "2x1" il sabato
    \item KPI: numero di ordini online, incremento delle visite, tasso di conversione
\end{itemize}

\subsubsection{Conversione di un potenziale cliente: esempio}
Supponiamo che una palestra investe in una campagna FAcebook/INstagram con video brevi che mostrano l'ambiente e i corsi. Pubblica contenuti su TikTok: brevi esercizi, consigli su allenamento e dieta e sponsorizza un evento gratuito all'aperto. 

Un utente vede il video, ne \'e incuiosito e visita il profilo Instagram della palestra, scoprendone cos\'i l'interesse.

Adesso il potenziale cliente ne vuole sapere di pi\'u. Il profilo Instagram \'e curato,, con testimonianze di clienti soddisfatti, video dei trainer, e storie in evidenza sui corsi, cos\'i come il sito web aggiornato con programma corsi e prezzi chiari. Il cliente visita il sito, legge le recensioni e si iscrive alla newsletter con un coupon del 10\%.

Essendosi iscritto, il cliente riceve via email il coupon e un invito a una prova gratuita con una breve descrizione dei corsi pi\'u adatti in base al suo interesse. A questo punto, il cliente prenota la lezione gratuita. 

Durante la prova, il personale \'e accogliente e professionale e a fine lezione, riceve una proposta limitata: "Iscriviti oggi, hai 1 mese gratis se porti un amico". Il risultato \'e che il cliente si iscrive e paga il primo mese. La conversione \'e avvenuta.

Il cliente dopo un mese diventa fedele e promuove il brand. Ogni settimana riceve via WhatsApp consigli personalizzati, orari corsi e promo solo per iscritti. Dopo 1 mese, la palestra chiede una recensione e offre un piccolo omaggio per ogni amico che si iscrive (programma referral).

\end{document}